% !TEX program = xelatex
% Lernplatt - by Can Siebert
% Kurs 71 (Dig 42) - Tag 15 - Lehrskript
% Nachhaltige Prozessveränderung führen
%
% Self-contained xelatex source. Kein biblatex, keine externen Bilder.

\documentclass[11pt,a4paper,oneside]{article}

% --- Encoding & Sprache --------------------------------------------------
\usepackage{polyglossia}
\setdefaultlanguage{german}

% --- Geometrie -----------------------------------------------------------
\usepackage[a4paper,top=2cm,bottom=2cm,outer=2.5cm,inner=2.5cm,bindingoffset=1.5cm]{geometry}

% --- Fonts ---------------------------------------------------------------
\usepackage{fontspec}
\defaultfontfeatures{Ligatures=TeX}

\IfFontExistsTF{Charter}{%
  \setmainfont{Charter}[%
    Numbers=OldStyle,%
    UprightFeatures   = {SmallCapsFeatures={Letters=SmallCaps}}%
  ]%
}{%
  \IfFontExistsTF{Bitstream Charter}{%
    \setmainfont{Bitstream Charter}%
  }{%
    \setmainfont{TeX Gyre Schola}%
  }%
}

\IfFontExistsTF{Source Sans 3}{%
  \setsansfont{Source Sans 3}%
}{%
  \IfFontExistsTF{Source Sans Pro}{%
    \setsansfont{Source Sans Pro}%
  }{%
    \setsansfont{TeX Gyre Heros}%
  }%
}

\newcommand{\caveatfontfallback}{\itshape}
\IfFontExistsTF{Caveat}{%
  \newfontfamily\caveatfont{Caveat}[Scale=1.15]%
}{%
  \newcommand\caveatfont{\caveatfontfallback}%
}

\setmonofont{Latin Modern Mono}

% --- Mikro-Typografie ----------------------------------------------------
\usepackage{microtype}
\linespread{1.15}

% --- Farben & Boxen ------------------------------------------------------
\usepackage{xcolor}
\definecolor{lpInk}{RGB}{20,28,38}
\definecolor{kurs71}{HTML}{9D4A1A}
\definecolor{lpAccent}{HTML}{9D4A1A}
\definecolor{lpAccentLight}{RGB}{248,232,218}
\definecolor{lpAmber}{RGB}{222,138,30}
\definecolor{lpAmberLight}{RGB}{255,243,217}
\definecolor{lpAmberBorder}{RGB}{196,118,18}
\definecolor{lpGray}{RGB}{96,104,114}
\definecolor{lpGrayLight}{RGB}{238,240,243}

\usepackage[most]{tcolorbox}
\tcbuselibrary{breakable,skins}

\newtcolorbox{eselsbox}[1][Eselsbrücke]{%
  enhanced, breakable,
  colback=lpAmberLight,
  colframe=lpAmberBorder,
  boxrule=0.6pt,
  arc=2pt,
  borderline={0.4pt}{0pt}{lpAmberBorder, dashed},
  left=10pt,right=10pt,top=8pt,bottom=8pt,
  fonttitle=\sffamily\bfseries\color{lpAmberBorder},
  coltitle=lpAmberBorder,
  title={#1},
  attach boxed title to top left={xshift=8pt,yshift=-8pt},
  boxed title style={size=small, colback=lpAmberLight, colframe=lpAmberBorder, boxrule=0.4pt, arc=1pt},
  top=14pt
}

\newtcolorbox{merkbox}[1][Merksatz]{%
  enhanced, breakable,
  colback=lpAccentLight,
  colframe=lpAccent,
  boxrule=0.6pt,
  arc=2pt,
  left=10pt,right=10pt,top=8pt,bottom=8pt,
  fonttitle=\sffamily\bfseries\color{white},
  coltitle=white,
  title={#1},
  attach boxed title to top left={xshift=8pt,yshift=-8pt},
  boxed title style={size=small, colback=lpAccent, colframe=lpAccent, boxrule=0.4pt, arc=1pt},
  top=14pt
}

% --- Listen --------------------------------------------------------------
\usepackage{enumitem}
\setlist{itemsep=2pt, topsep=4pt, parsep=0pt}
\setlist[itemize,1]{label={\textbullet},leftmargin=1.6em}
\setlist[itemize,2]{label={\textendash},leftmargin=1.4em}
\setlist[enumerate,1]{label=\arabic*., leftmargin=1.8em}

% --- Tabellen ------------------------------------------------------------
\usepackage{tabularx}
\usepackage{array}
\usepackage{booktabs}
\usepackage{longtable}

% --- Kopf-/Fusszeilen ----------------------------------------------------
\usepackage{fancyhdr}
\usepackage{lastpage}
\pagestyle{fancy}
\fancyhf{}
\renewcommand{\headrulewidth}{0.3pt}
\renewcommand{\footrulewidth}{0pt}
\fancyhead[L]{\footnotesize\sffamily\color{lpGray}Lernplatt \textperiodcentered{} Kurs 71 \textperiodcentered{} Tag 15}
\fancyhead[R]{\footnotesize\sffamily\color{lpGray}Change \textperiodcentered{} Stakeholder \textperiodcentered{} Sustainment}
\fancyfoot[L]{\footnotesize\sffamily\color{lpGray}\textcopyright{} Can Siebert}
\fancyfoot[R]{\footnotesize\sffamily\color{lpGray}\thepage{} / \pageref{LastPage}}

\fancypagestyle{plain}{%
  \fancyhf{}%
  \renewcommand{\headrulewidth}{0pt}%
  \fancyfoot[C]{\footnotesize\sffamily\color{lpGray}\thepage{} / \pageref{LastPage}}%
}

% --- Section-Styling -----------------------------------------------------
\usepackage[explicit]{titlesec}

\titleformat{\section}[block]
  {\sffamily\Large\bfseries\color{lpAccent}}
  {\thesection.}{0.6em}{#1}
  [{\vspace{2pt}\color{lpAccent}\titlerule[0.6pt]}]

\titleformat{\subsection}[block]
  {\sffamily\large\bfseries\color{lpInk}}
  {\thesubsection}{0.6em}{#1}

\titleformat{\subsubsection}[block]
  {\sffamily\normalsize\bfseries\color{lpInk}}
  {}{0pt}{#1}

\titlespacing*{\section}{0pt}{14pt}{8pt}
\titlespacing*{\subsection}{0pt}{10pt}{4pt}
\titlespacing*{\subsubsection}{0pt}{8pt}{2pt}

% --- Hyperlinks ----------------------------------------------------------
\usepackage[hidelinks,unicode]{hyperref}

% --- TOC -----------------------------------------------------------------
\renewcommand{\contentsname}{\sffamily\Large\bfseries\color{lpAccent}Inhalt}

% --- Helfer --------------------------------------------------------------
\newcommand{\akzent}[1]{{\caveatfont\color{lpAmber}#1}}
\newcommand{\fachbegriff}[1]{\textbf{#1}}
\newcommand{\quelle}[1]{{\footnotesize\color{lpGray}(#1)}}

% =========================================================================
\begin{document}

% --- COVER ---------------------------------------------------------------
\begin{titlepage}
  \pagestyle{empty}
  \vspace*{1.5cm}
  \noindent{\sffamily\footnotesize\color{lpGray}LERNPLATT \textperiodcentered{} BY CAN SIEBERT}\\[2pt]
  \noindent{\color{lpAccent}\rule{\linewidth}{1.2pt}}\\[4pt]
  \noindent{\sffamily\footnotesize\color{lpGray}MODUL 7 \textperiodcentered{} TAG 15 VON 16 \textperiodcentered{} KURS 71 (Dig 42) \textperiodcentered{} 18.\,Juni 2026}

  \vspace{3cm}
  {\sffamily\fontsize{30}{34}\selectfont\bfseries\color{lpInk}Prozess-\\veränderung\\nachhaltig führen\par}

  \vspace{0.7cm}
  {\sffamily\large\color{lpGray}Rollenklärung, Stakeholder, Widerstand produktiv machen --\\Befähigen und Verankern mit messbarer Wirksamkeit.\par}

  \vspace{1.5cm}
  {\caveatfont\LARGE\color{lpAmber}Lehrskript -- Tag 15 von 16}\par

  \vfill

  \noindent\begin{minipage}{0.55\linewidth}
    {\sffamily\footnotesize\color{lpGray}Lernpfad:}\\
    {\sffamily\small\color{lpInk}Wissen \textrightarrow{} Anwendung \textrightarrow{} Management}\\[10pt]
    {\sffamily\footnotesize\color{lpGray}Bearbeitungsdauer:}\\
    {\sffamily\small\color{lpInk}1 Kurstag}
  \end{minipage}\hfill
  \begin{minipage}{0.4\linewidth}
    \raggedleft
    {\sffamily\footnotesize\color{lpGray}Dozent:}\\
    {\sffamily\small\color{lpInk}Can Siebert}\\[10pt]
    {\sffamily\footnotesize\color{lpGray}Format:}\\
    {\sffamily\small\color{lpInk}Theorie \textperiodcentered{} Fallstudie \textperiodcentered{} Coaching}
  \end{minipage}

  \vspace{0.5cm}
  \noindent{\color{lpAccent}\rule{\linewidth}{0.6pt}}
\end{titlepage}

% --- LERNZIELE -----------------------------------------------------------
\section*{Lernziele dieses Tages}
\addcontentsline{toc}{section}{Lernziele dieses Tages}
\thispagestyle{plain}

\noindent Prozessveränderungen scheitern selten an der Methode -- sie scheitern am Menschen, an unklaren Rollen und an fehlender Verankerung. Heute lernen Sie, Veränderungen \emph{nachhaltig zu führen}: Rollen klären, Stakeholder strukturiert einbinden, Widerstand als Information lesen, Befähigung gezielt ermöglichen und Wirkung messbar verankern. Aus ``Change als Projekt'' wird ``Change als Betriebssystem''.

\subsection*{L1 -- Wissen (Reproduktion und Einordnung)}
\begin{itemize}
  \item Change-Leadership-Rollen kennen: Sponsor, Change Lead, Process Owner, Line Manager, Champions.
  \item Erfolgsfaktoren benennen: klares ``Warum'', Beteiligung, Quick Wins, psychologische Sicherheit, sichtbare Führung.
  \item Widerstand als Information verstehen: Interessen, Identität, Überlastung, Vertrauen -- nicht als ``Sturheit''.
  \item Sustainment-Mechanik kennen: Routinen, KPIs, Reviews, Adoption-Messung.
\end{itemize}

\subsection*{L2 -- Anwendung (Transfer auf konkrete Situationen)}
\begin{itemize}
  \item Eine Change-Architektur entwerfen (Rollen, Entscheidungswege, Rituale).
  \item Eine Stakeholderstrategie 2$\times$2 (Einfluss / Betroffenheit) mit Einbindungsstufen erstellen.
  \item Ein Kommunikationspaket mit Kernbotschaften und Feedback-Mechanismen formulieren.
  \item Ein Sustainment-Design mit Outcome-, Prozess- und Adoption-KPIs verankern.
\end{itemize}

\subsection*{L3 -- Management (Entscheidung und Verantwortung)}
\begin{itemize}
  \item Transparenz-Regeln für KPIs entscheiden (was wird sichtbar, was nicht, warum).
  \item Stopp- oder Schwenk-Kriterien für Change-Vorhaben festlegen.
  \item WIP-Limits für Change-Portfolio definieren -- Überlastung vermeiden.
  \item Trade-off-Protokolle führen (Speed vs.\ Compliance, Transparenz vs.\ Angst).
\end{itemize}

\begin{merkbox}[Roter Faden]
Change ist nicht das, was im Folien-Plan steht -- Change ist das, was nach 6 Monaten noch gilt. Die Disziplin: nicht ``managen'', sondern \emph{führen}. Klare Rollen, ehrliche Kommunikation, Widerstand produktiv machen, befähigen, verankern. Wer das nicht durchhält, baut wiederholt dieselben Projekte ohne dauerhaften Effekt. \quelle{Kotter, 1996; Lewin, 1951; Schein, 2017}
\end{merkbox}

\clearpage

% --- 1. EINSTIEG ---------------------------------------------------------
\section{Einstieg: Vom Projekt zum Betriebssystem}

Eine Prozessveränderung verläuft selten linear. Sie hat Phasen, in denen alles möglich erscheint, und Phasen, in denen Widerstand und Müdigkeit das Tempo brechen. Diese Dynamik ist nicht ``Pech'' -- sie ist die Regel. Wer sie kennt, kann sie führen.

\subsection{Drei klassische Modelle}

\begin{enumerate}
  \item \fachbegriff{Lewin (1951)}: \emph{Unfreeze -- Change -- Refreeze.} Das Modell hat sich gehalten, weil es eine wichtige Wahrheit fasst: Veränderung beginnt mit Auftauen alter Muster und endet mit Verankern der neuen. Ohne Refreeze ist alles temporär. \quelle{Lewin, 1951}
  \item \fachbegriff{Kotter (1996)}: \emph{Acht Schritte} (Sense of Urgency $\to$ Coalition $\to$ Vision $\to$ Communicate $\to$ Empower $\to$ Quick Wins $\to$ Sustain $\to$ Institutionalize). Pragmatischer Leitfaden, der in vielen Organisationen Sprache schafft. \quelle{Kotter, 1996}
  \item \fachbegriff{Schein -- Kulturmodell}: drei Schichten: \emph{Artefakte} (sichtbar), \emph{Werte} (kommuniziert), \emph{Grundannahmen} (unbewusst). Nachhaltige Veränderung greift in die zweite und dritte Schicht ein, nicht nur in die erste. \quelle{Schein, 2017}
\end{enumerate}

\subsection{Change als Betriebssystem, nicht als Projekt}

Ein Projekt hat ein Start- und ein Enddatum. Ein Betriebssystem läuft. Drei Indikatoren, dass Change zum Betriebssystem geworden ist:

\begin{itemize}
  \item Es gibt klare Rollen für Sustainment (nicht nur für die Einführung).
  \item Es gibt Routinen (Standups, Reviews, Improvement Boards), die nach dem ``Go-Live'' weiterlaufen.
  \item Es gibt KPIs, die kontinuierlich gemessen werden -- nicht nur in der Projektphase.
\end{itemize}

\begin{eselsbox}[Eselsbrücke: \akzent{U-C-R}]
Lewin in drei Buchstaben:\\[2pt]
\textbf{U}nfreeze -- alte Muster auftauen, Dringlichkeit schaffen\\
\textbf{C}hange -- neue Praktiken einführen, lernen lassen\\
\textbf{R}efreeze -- verankern, sodass es nicht zurückkippt\\[2pt]
\emph{Wer das Refreeze auslässt, beginnt das Projekt in 12 Monaten von vorn.}
\end{eselsbox}

% --- 2. ROLLEN -----------------------------------------------------------
\section{Change-Leadership-Rollen: wer macht was?}

In jedem Change-Vorhaben treten dieselben Rollen auf -- mit unterschiedlichen Verantwortungen. Klare Rollen sind die Basis für alles Weitere.

\subsection{Die fünf Kernrollen}

\begin{itemize}
  \item \fachbegriff{Sponsor.} Gibt Richtung, Ressourcen, schützt die Veränderung gegen Widerstände auf höherer Ebene, entscheidet bei Konflikten. Ohne Sponsor stirbt der Change beim ersten ernsten Widerstand.
  \item \fachbegriff{Change Lead.} Orchestriert Kommunikation, Stakeholder, Enablement und Messung. Häufige Aufgabe: Übersetzer zwischen Sponsor, Process Owner, Linie und Champions.
  \item \fachbegriff{Process Owner.} Fachliche Verantwortung. Setzt Standards, definiert KPIs, sichert das Sustainment im Alltag.
  \item \fachbegriff{Line Manager.} Übersetzt die Veränderung in den Alltag des Teams. Entfernt Hindernisse, führt über Verhalten, gibt Feedback.
  \item \fachbegriff{Champions / Multiplikatoren.} Mitarbeitende mit Nähe zur Praxis, informellem Vertrauen, Bereitschaft, Verbesserungen zu testen und Feedback zu sammeln.
\end{itemize}

\subsection{Rollen-Verantwortungen je drei Punkten}

\begin{tabularx}{\linewidth}{@{}>{\bfseries}lX@{}}
\toprule
Rolle & Drei Kernverantwortungen \\
\midrule
Sponsor & Ressourcen sichern; bei Konflikten entscheiden; Veränderung schützen \\
Change Lead & Stakeholder orchestrieren; Kommunikation gestalten; Wirkung messen \\
Process Owner & Standards setzen; KPIs definieren; Sustainment sichern \\
Line Manager & Hindernisse entfernen; Verhalten coachen; Feedback einsammeln \\
Champions & Verbesserung testen; Feedback bündeln; Vertrauen brücken \\
\bottomrule
\end{tabularx}

\subsection{Risikoliste typischer Change-Vorhaben}

\begin{itemize}
  \item Akzeptanzproblem -- ``Wieder ein Projekt''.
  \item Überlastung -- parallele Initiativen verdrängen einander.
  \item Zielkonflikt -- Speed vs.\ Compliance, Transparenz vs.\ Autonomie.
  \item Fehlende Daten -- Baseline unklar, Wirkung kaum messbar.
  \item Schattenprozesse -- offizieller Prozess wird umgangen.
  \item Fehlende Sponsor-Klarheit -- niemand entscheidet bei Konflikten.
  \item Trainings ohne Praxisbezug -- Wissen ohne Verhaltenstransfer.
  \item Mitarbeiterverlust -- Schlüsselpersonen verlassen das Unternehmen.
\end{itemize}

\subsection{Risiko-Hypothesen testen}

Senior-Disziplin: nicht alle Risiken werden ``vermieden'', manche werden \emph{getestet}. In den ersten zwei Wochen werden zwei Hypothesen aktiv überprüft -- mit Mini-Experimenten:

\begin{itemize}
  \item Hypothese 1: ``Operations hat keine Zeit für Schulungen.'' Test: Erste Pilot-Schulung von 30 Minuten anbieten, Teilnahmequote messen.
  \item Hypothese 2: ``Fachbereich fürchtet KPI-Kontrolle.'' Test: Workshop zu KPI-Designprinzipien ($\to$ ``für Lernen, nicht für Bestrafung''), Feedback einholen.
\end{itemize}

% --- 3. STAKEHOLDER ------------------------------------------------------
\section{Stakeholder-Strategie: wer wird wie eingebunden?}

Eine gute Stakeholderstrategie unterscheidet sich von einer schlechten in einem Punkt: sie macht aus Stakeholdern \emph{Beteiligte}, nicht ``Betroffene''. Wer als Beteiligter angesprochen wird, hat eine Bringschuld; wer als Betroffener behandelt wird, eine Holschuld nach Information.

\subsection{Die 2$\times$2-Matrix: Einfluss vs.\ Betroffenheit}

Eine pragmatische Stakeholder-Klassifikation:

\begin{tabularx}{\linewidth}{@{}l XX@{}}
\toprule
& \textbf{Niedrige Betroffenheit} & \textbf{Hohe Betroffenheit} \\
\midrule
\textbf{Hoher Einfluss} & Konsultieren & Mitentscheiden \\
\textbf{Niedriger Einfluss} & Informieren & Verantworten \\
\bottomrule
\end{tabularx}

\medskip

Die vier Einbindungsstufen:

\begin{itemize}
  \item \fachbegriff{Informieren.} Status, Termine, Auswirkungen werden mitgeteilt. Ein Kanal genügt, Frequenz angemessen.
  \item \fachbegriff{Konsultieren.} Vor Entscheidungen werden diese Stakeholder gefragt. Nicht stimmberechtigt, aber stimmhabend.
  \item \fachbegriff{Mitentscheiden.} Stakeholder ist Teil der Entscheidung. Voting oder Konsens je nach Thema.
  \item \fachbegriff{Verantworten.} Stakeholder trägt die Konsequenz; ist Teil von Design und Umsetzung.
\end{itemize}

\subsection{Acht typische Stakeholder in Prozessveränderungen}

\begin{itemize}
  \item Operations (Tagesbetrieb).
  \item Fachbereiche (fachliche Verantwortung).
  \item IT (technische Umsetzung).
  \item Compliance (regulatorische Anforderungen).
  \item Controlling (KPI-Logik, Wirtschaftlichkeit).
  \item Customer Service (Kundensicht).
  \item Betriebsrat (Mitbestimmung, wenn relevant).
  \item Key Customer (externe Wirkung).
\end{itemize}

\subsection{Drei klassische Blockaden}

\begin{itemize}
  \item \textbf{Team A (Operations).} ``Schon wieder ein Projekt, wir haben keine Zeit.'' Antwort: WIP-Limit für Change, Entlastung sichtbar machen, Quick Wins zuerst.
  \item \textbf{Team B (Fachbereich).} ``KPI-Steuerung = Kontrolle, wir verlieren Autonomie.'' Antwort: Designprinzipien transparent machen, KPIs für Lernen statt Bestrafung, Beteiligung am KPI-Design.
  \item \textbf{Team C (IT).} ``Ohne klare Anforderungen wird das nichts.'' Antwort: gemeinsame Klärung der Datenquellen, frühe Architektur-Skizze, IT als Mitentscheider.
\end{itemize}

\begin{merkbox}[Widerstand als Information]
Widerstand ist selten ``Sturheit''. Er ist meistens \emph{Information}: über Interessen, Identität, Überlastung, fehlendes Vertrauen, ungelöste Trade-offs. Die Disziplin: zuhören, bevor man überzeugt. Eine Veränderung, die alle Argumente versteht und keines übergeht, ist robuster als eine, die ``durchgedrückt'' wird.
\end{merkbox}

% --- 4. KOMMUNIKATION ----------------------------------------------------
\section{Kommunikationsarchitektur: Botschaft, Kanal, Feedback}

Eine Kommunikationsstrategie ist nicht ``viele E-Mails'', sondern eine \emph{Architektur} aus Zielgruppe, Botschaft, Kanal, Frequenz, Feedback. \quelle{Kotter, 1996, Kap.~6}

\subsection{Drei Botschaftstypen}

\begin{enumerate}
  \item \fachbegriff{Was?} Was ändert sich konkret? Welche Prozesse, welche Werkzeuge, welche Termine?
  \item \fachbegriff{Warum?} Warum jetzt, warum so? Was passiert, wenn wir nicht verändern? Das ``Sense of Urgency'' nach Kotter.
  \item \fachbegriff{Was ändert sich für mich?} Die individuelle Sicht. Ohne diese Antwort bleibt jede Botschaft abstrakt.
\end{enumerate}

\subsection{Kanal-Mix}

Ein Kanal reicht selten. Pragmatischer Mix:

\begin{itemize}
  \item \textbf{Sponsor-Kommunikation} -- bei Meilensteinen, in eigener Sprache.
  \item \textbf{Town Hall} -- Q\&A live, mit dem Mut, schwierige Fragen zu beantworten.
  \item \textbf{Team-Meetings} -- Line Manager übersetzen für ihren Bereich.
  \item \textbf{Schriftliche Kanäle} -- Newsletter, FAQ, Intranet-Seite mit aktuellem Stand.
  \item \textbf{Feedback-Kanäle} -- Pulse Survey, anonymes Signal, offene Sprechstunde.
\end{itemize}

\subsection{Feedback-Mechanismen}

Kommunikation ohne Feedback ist eine Predigt. Drei Mechanismen:

\begin{itemize}
  \item \textbf{Pulse Survey} (kurz, häufig). 3--5 Fragen, monatlich.
  \item \textbf{Anonymes Signal} (z.\,B.\ ein einfaches Formular für ``Sorgen / Hindernisse''). Vertrauliche Rückläufe.
  \item \textbf{Offene Q\&A} (Sprechstunde, Town Hall). Live, mit Bereitschaft, ``Ich weiß es nicht'' zu sagen.
\end{itemize}

\begin{eselsbox}[Eselsbrücke: \akzent{W-W-M}]
Drei Kernbotschaften, die in jedem Change-Vorhaben beantwortet sein müssen:\\[2pt]
\textbf{W}as ändert sich konkret?\\
\textbf{W}arum jetzt, warum so?\\
\textbf{M}ich -- was ändert sich für mich?\\[2pt]
\emph{Wer die dritte Frage übergeht, redet an den Menschen vorbei.}
\end{eselsbox}

% --- 5. SUSTAINMENT ------------------------------------------------------
\section{Sustainment: Verankern statt verpuffen}

Die meisten Veränderungen scheitern nicht am Start, sondern am Sustainment. Nach den ersten Erfolgen bricht die Aufmerksamkeit ab, neue Themen drängen sich vor, alte Muster kehren zurück.

\subsection{Vier Säulen des Sustainments}

\begin{enumerate}
  \item \fachbegriff{Standards.} Was ist die neue ``Definition of Done''? Was wird wie protokolliert?
  \item \fachbegriff{Routinen.} Regelmäßige Rituale (Weekly Standup, Monthly Improvement Board) tragen die Veränderung über Zeit.
  \item \fachbegriff{Messung.} KPIs werden kontinuierlich erhoben; Abweichungen werden besprochen.
  \item \fachbegriff{Coaching.} Line Manager und Champions coachen weiter; nicht ``Schulung am Anfang, dann fertig''.
\end{enumerate}

\subsection{Fünf-Dimensions-KPIs}

Wirksamkeit lässt sich in fünf Dimensionen messen:

\begin{itemize}
  \item \fachbegriff{Outcome.} Hat sich das Geschäftsergebnis verbessert? (Kunde, Qualität, Kosten.)
  \item \fachbegriff{Prozess.} DLZ kürzer, Rework geringer, SLA-Quote besser?
  \item \fachbegriff{Adoption.} Wieviele Fälle laufen im neuen Prozess? Wieviele umgehen ihn?
  \item \fachbegriff{Capability.} Wie viele Mitarbeitende haben Training und Coaching abgeschlossen?
  \item \fachbegriff{Culture.} Psychologische Sicherheit, Ownership, Lernbereitschaft (Pulse Survey).
\end{itemize}

\subsection{Definition of Done für Prozessänderungen}

Eine pragmatische DoD:

\begin{itemize}
  \item Modell aktuell und freigegeben.
  \item Owner benannt und akzeptiert.
  \item KPIs adjustiert, im Dashboard sichtbar.
  \item Training durchgeführt für relevante Rollen.
  \item Monitoring aktiv.
  \item Erste Iteration der Improvement-Schleife durchlaufen.
\end{itemize}

\subsection{Anti-Gaming bei KPIs}

KPIs ohne Schutz vor Manipulation erzeugen oft das Gegenteil dessen, was sie messen sollen:

\begin{itemize}
  \item KPI ``Tickets pro Stunde'' $\to$ Anreiz, Tickets zu splitten.
  \item KPI ``First-Time-Right'' $\to$ Anreiz, schwierige Fälle zu vermeiden.
  \item KPI ``Antwortzeit'' $\to$ Anreiz, leere Antworten zu senden.
\end{itemize}

Schutz: KPIs koppeln (z.\,B.\ ``Tickets pro Stunde'' nur gültig, wenn Reklamationsquote unter X bleibt). Und: KPI-Transparenz \emph{auf Team-Ebene}, nicht auf Personen-Ebene -- außer bei begründetem Einzelfall.

\begin{merkbox}[KPIs für Lernen, nicht für Bestrafung]
Die wichtigste Designentscheidung in KPI-Systemen: was tun wir, wenn ein KPI nicht erreicht wird? Wenn die Antwort ``Konsequenz für die Person'' ist, wird der KPI manipuliert. Wenn die Antwort ``Lernen, was hindert'' ist, wird er ein Steuerungsinstrument. Diese Entscheidung muss vor der KPI-Einführung gefallen sein.
\end{merkbox}

% --- 6. FALLSTUDIE SERVICECO ---------------------------------------------
\section{Fallstudie: ServiceCo -- E2E-Prozessveränderung}

\emph{ServiceCo} hat mehrere Standorte und führt einen workflowgestützten End-to-End-Prozess inkl.\ KPI-Dashboard und Abweichungsmanagement ein. Ziel: $-$25\,\% Durchlaufzeit, $-$30\,\% Rework in 12 Wochen.

Restriktionen: Zeitdruck, parallele Projekte, unvollständige Daten, Widerstand gegen KPI-Transparenz, Budget begrenzt.

\subsection{Change-Architektur}

\begin{itemize}
  \item \textbf{Sponsor.} Bereichsleiter Operations. Entscheidet Zielkonflikte und Ressourcen.
  \item \textbf{Process Owner.} Senior aus Operations. Accountable für Standards und KPIs.
  \item \textbf{Change Lead.} Interne Beraterin. Orchestriert Kommunikation und Enablement.
  \item \textbf{Line Manager.} An jedem Standort, führen Verhalten.
  \item \textbf{Champions.} 2--3 pro Standort, sammeln Feedback und testen Verbesserungen.
\end{itemize}

Rituale:

\begin{itemize}
  \item \textbf{Weekly Process Standup} (30\,min). Standort-Lead + Champions. Kennzahlen, Hindernisse.
  \item \textbf{Monthly Improvement Board} (60\,min). Standort-übergreifend. Trade-offs, Standardanpassungen.
  \item \textbf{Quarterly Review} (je nach Kontext). Sponsor + Process Owner + ggf.\ Compliance.
\end{itemize}

\subsection{Stakeholderstrategie}

\begin{itemize}
  \item Operations: Verantworten (Tagesbetrieb).
  \item Fachbereich: Mitentscheiden (KPI-Design).
  \item IT: Mitentscheiden (Controls, Datenfluss).
  \item Compliance: Mitentscheiden (Audit-Anforderungen).
  \item Controlling: Konsultieren (KPI-Logik).
  \item Customer Service: Konsultieren (Sicht des Kunden).
  \item Betriebsrat: Konsultieren / Mitentscheiden (Datenschutz, Personenbezug).
  \item Key Customer: Informieren (Veränderungen im Service).
\end{itemize}

\subsection{Kommunikationspaket}

\begin{itemize}
  \item \textbf{Kick-off-Botschaft.} ``Wir reduzieren Rework, schützen Qualität, geben Teams Klarheit. KPIs sind zur Verbesserung, nicht zur Bestrafung.''
  \item \textbf{Drei Kernbotschaften.} (1) Was sich konkret ändert; (2) Warum jetzt; (3) Was sich für mich ändert (rollenspezifisch).
  \item \textbf{Zwei Kanäle.} Town Hall am Anfang + monatlicher Newsletter mit Fortschritt.
  \item \textbf{Feedback.} Offene Q\&A + anonymer Pulse Survey monatlich.
\end{itemize}

\subsection{Sustainment-Design}

\begin{itemize}
  \item \textbf{KPI-Set.}
    \begin{itemize}
      \item Outcome: SLA-Quote, Kundenzufriedenheit.
      \item Prozess: DLZ, Rework.
      \item Adoption: Anteil Fälle im Workflow, Exception-Rate.
      \item Capability: Trainingsabschluss, Coaching-Check.
    \end{itemize}
  \item \textbf{Review-Board.} Monthly Improvement Board, mit Protokoll und Aktionen.
  \item \textbf{Trainings- und Coaching-Plan.} Drei Stufen: Operator-Schulung, Champion-Workshop, Coaching für Line Manager.
  \item \textbf{Definition of Done.} Modell aktuell, Owner benannt, KPIs angepasst, Training durchgeführt, Monitoring aktiv.
\end{itemize}

\subsection{Entscheidungen unter Unsicherheit}

\begin{enumerate}
  \item \textbf{Speed vs.\ Compliance.} MVP mit Ausnahmeprozess + Audit Trail; nicht ``erst alle Kontrollen perfekt, dann starten''.
  \item \textbf{Transparenz-Regel.} Team-Performance sichtbar; individuelle Personenmetriken nur bei begründetem Incident-Case (Guardrail).
\end{enumerate}

Trade-off-Protokoll: jede dieser Entscheidungen wird im Improvement Board protokolliert -- inklusive verworfener Alternativen und Frühwarnsignal (``Wenn KPI-Akzeptanz unter X fällt, reviewen wir die Transparenz-Regel.'').

% --- 7. COACHING ---------------------------------------------------------
\section{Coaching: Führung = Verhalten, Entscheidungen, Schutz}

Change-Leadership ist nicht ``mehr Folien''. Es ist eine Kombination aus drei Disziplinen:

\subsection{Verhalten}

\begin{itemize}
  \item Vorleben statt verkünden. Ein Sponsor, der die neuen Routinen ignoriert, hebelt sie sofort aus.
  \item Sichtbar in den Rituals. Erscheinen, fragen, zuhören.
  \item ``Was tust du ab morgen anders?'' -- die Frage, die Verhaltenstransfer sichert.
\end{itemize}

\subsection{Entscheidungen}

\begin{itemize}
  \item Zielkonflikte explizit entscheiden. Niemand kann gleichzeitig Speed, Qualität, Compliance und niedrige Kosten optimieren.
  \item Klare Stop-Bedingungen. Wann wird ein Change-Vorhaben gestoppt oder neu geschnitten?
  \item Trade-off-Protokoll. Jede schwierige Entscheidung mit verworfenen Alternativen und Frühwarnsignal.
\end{itemize}

\subsection{Schutz}

\begin{itemize}
  \item Teams vor parallelen Initiativen schützen (WIP-Limit für Change).
  \item Psychologische Sicherheit -- wer Fehler meldet, wird nicht bestraft.
  \item Ressourcen sichern, wenn der Sponsor an anderer Stelle reduzieren will.
\end{itemize}

\subsection{Coaching-Interventionen}

\begin{itemize}
  \item \textbf{Aktives Zuhören.} ``Was genau bereitet Ihnen Sorge?''
  \item \textbf{Reframing.} ``Sehe ich richtig, dass Sie auf ein Risiko hinweisen -- nämlich X?''
  \item \textbf{Konfliktmoderation.} Trade-offs sichtbar machen statt zu schlichten.
  \item \textbf{Commitment.} ``Was tun Sie ab morgen anders?''
\end{itemize}

% --- 8. MANAGEMENT-PERSPEKTIVE -------------------------------------------
\section{Management-Perspektive: Change als Lernsystem}

Auf Wissens-Ebene ist Change ein Modell. Auf Anwendungs-Ebene ist es ein Projekt. Auf Management-Ebene ist es ein \emph{Lernsystem} -- mit Portfolio, Cadence, Feedback und Korrektur.

\subsection{Top-10-Change-Entscheidungen}

\begin{enumerate}
  \item Welche Initiativen werden parallel zugelassen (WIP-Limit)?
  \item Welche Ressourcen sind reserviert (Sponsor-Disziplin)?
  \item Welche KPI-Transparenz wird zugelassen (Team / Person)?
  \item Welcher Ausnahmeprozess gilt?
  \item Welche Eskalationspfade existieren?
  \item Welche Kommunikationsregeln (Frequenz, Kanäle)?
  \item Welche Trainings-Investitionen sind verbindlich?
  \item Welche Stopp-Kriterien werden definiert?
  \item Welche Rituale sind verbindlich (Standup, Improvement Board)?
  \item Welche Sustainment-DoD gilt?
\end{enumerate}

\subsection{Portfolio-Steuerung für Change}

\begin{itemize}
  \item \textbf{WIP-Limit.} Eine Organisation hält eine bestimmte Zahl gleichzeitiger Veränderungen aus. Wer mehr startet, riskiert, dass alle scheitern.
  \item \textbf{Priorisierung.} Welche Initiativen sind ``must do'', welche ``nice to do''?
  \item \textbf{Entlastungs-Mechanismus.} Wenn eine neue Initiative startet, was wird parallel pausiert?
\end{itemize}

\subsection{Truth Signals}

Wie erkennt man, ob eine Veränderung wirklich angekommen ist -- nicht nur ``offiziell läuft''?

\begin{itemize}
  \item \textbf{Adoption-Rate.} Wieviele Fälle laufen tatsächlich im neuen Prozess?
  \item \textbf{Schattenprozess-Indikatoren.} Werden Workarounds (Excel, Mail) immer noch genutzt?
  \item \textbf{Pulse Survey.} Wie schätzen Mitarbeitende selbst den Stand ein?
  \item \textbf{Outcome-Verbesserung.} Hat sich das gemessene Geschäftsergebnis bewegt?
\end{itemize}

\subsection{Risikomodell}

\begin{itemize}
  \item \textbf{Widerstandstypen.} Identität, Überlastung, Vertrauen, Interessen.
  \item \textbf{Überlastungsindikatoren.} Krankenstand, Pulse Survey, sinkende Adoption-Rate.
  \item \textbf{KPI-Gaming-Erkennung.} Plötzliche, unplausible Verbesserung; Diskrepanz zwischen KPI und Kundenfeedback.
  \item \textbf{Schatten-Prozesse.} Anstieg von Workarounds, sinkende Workflow-Quote.
\end{itemize}

\begin{merkbox}[Schluss-Merksatz für Tag 15]
Change ist ein \textbf{Lernsystem}: klare Rollen, ehrliche Kommunikation, ehrliche Messung, Widerstand als Information, kontinuierliche Verbesserung. Wer Change ``managen'' will wie ein Projekt, baut wiederkehrende Initiativen ohne Verankerung. Wer Change \emph{führt} wie ein Betriebssystem, schafft eine Organisation, die sich verändern \emph{kann}.
\end{merkbox}

% --- 9. TAKE-AWAYS -------------------------------------------------------
\section{Take-aways aus Tag 15}

\begin{enumerate}
  \item \textbf{Lewin, Kotter, Schein} sind die drei Grundmodelle -- Unfreeze/Change/Refreeze, 8 Steps, Kulturmodell.
  \item \textbf{Fünf Kernrollen} klären: Sponsor, Change Lead, Process Owner, Line Manager, Champions.
  \item \textbf{Widerstand ist Information}, keine Sturheit. Erst zuhören, dann gestalten.
  \item \textbf{Stakeholder 2$\times$2}: Informieren / Konsultieren / Mitentscheiden / Verantworten.
  \item \textbf{Drei Kernbotschaften} pro Change-Vorhaben: Was, Warum, Was für mich.
  \item \textbf{Sustainment in vier Säulen}: Standards, Routinen, Messung, Coaching.
  \item \textbf{KPIs für Lernen, nicht für Bestrafung}; mit Kopplungen gegen Gaming.
  \item \textbf{Change ist Lernsystem}, kein Projekt. WIP-Limits und Truth Signals als Steuerungsinstrumente.
\end{enumerate}

% --- 10. LITERATUR -------------------------------------------------------
\clearpage
\section{Literatur}

Im Skript verwendete und für die Vertiefung empfohlene Quellen. Zitierstil: APA-7.

\begin{description}[style=nextline,leftmargin=18pt,labelindent=0pt,itemsep=4pt]
  \item[Edmondson, A.\,C. (1999).] Psychological safety and learning behavior in work teams. \emph{Administrative Science Quarterly, 44}(2), 350--383. \href{https://doi.org/10.2307/2666999}{https://doi.org/10.2307/2666999}
  \item[Kotter, J.\,P. (1996).] \emph{Leading change.} Harvard Business School Press.
  \item[Kotter, J.\,P. (2012).] \emph{Leading change: With a new preface.} Harvard Business Review Press.
  \item[Lewin, K. (1951).] \emph{Field theory in social science: Selected theoretical papers.} Harper \& Row.
  \item[Project Management Institute. (2021).] \emph{A guide to the Project Management Body of Knowledge (PMBOK\textregistered{} Guide)} (7th ed.). Project Management Institute.
  \item[Rogers, E.\,M. (2003).] \emph{Diffusion of innovations} (5th ed.). Free Press.
  \item[Schein, E.\,H. (2017).] \emph{Organizational culture and leadership} (5th ed.). Jossey-Bass.
  \item[Westerman, G., Bonnet, D., \& McAfee, A. (2014).] \emph{Leading digital: Turning technology into business transformation.} Harvard Business Review Press.
\end{description}

% --- 11. GLOSSAR ---------------------------------------------------------
\clearpage
\section{Glossar}

Die wichtigsten Begriffe des Tages -- kurz definiert, in alphabetischer Reihenfolge.

\begin{description}[style=nextline,leftmargin=0pt,labelindent=0pt,itemsep=4pt]
  \item[Adoption-Rate] Anteil der Fälle, die tatsächlich im neuen Prozess laufen. Truth-Signal-Indikator.
  \item[Anti-Gaming (KPI)] Designprinzip: KPIs koppeln, sodass Manipulation an einer Stelle eine andere verschlechtert.
  \item[Champion] Mitarbeitender mit Nähe zur Praxis und informellem Vertrauen, der Verbesserungen testet und Feedback sammelt.
  \item[Change Lead] Person, die Kommunikation, Stakeholder und Enablement in einer Veränderung orchestriert.
  \item[Change-Architektur] Sammlung der Rollen, Entscheidungswege und Rituale einer Veränderung.
  \item[Definition of Done (Change)] Pragmatische Kriterienliste, wann eine Prozessänderung wirklich abgeschlossen ist.
  \item[Improvement Board] Regelmäßiges Gremium (z.\,B.\ monatlich), das Trade-offs und Standardanpassungen entscheidet.
  \item[Kotter (8 Steps)] Modell zur Führung von Veränderungen in acht Phasen. \quelle{Kotter, 1996}
  \item[Lewin (Unfreeze/Change/Refreeze)] Dreischrittmodell der Veränderung. \quelle{Lewin, 1951}
  \item[Line Manager] Führungskraft, die die Veränderung in den Alltag des Teams übersetzt.
  \item[Process Owner] Fachliche Verantwortung für Standards, KPIs und Sustainment.
  \item[Pulse Survey] Kurze, häufige Befragung zu Stand und Stimmung; 3--5 Fragen, monatlich.
  \item[Psychologische Sicherheit] Klima, in dem Mitarbeitende Risiken eingehen und Fehler melden können, ohne bestraft zu werden. \quelle{Edmondson, 1999}
  \item[Refreeze] Verankerung neuer Praktiken (Standards, Routinen, KPIs), sodass Rückfall verhindert wird.
  \item[Rituale] Verbindliche, wiederkehrende Termine wie Standup, Review, Improvement Board.
  \item[Schein-Kulturmodell] Dreischicht-Modell der Organisationskultur: Artefakte, Werte, Grundannahmen. \quelle{Schein, 2017}
  \item[Sponsor] Senior-Führungskraft, die Richtung, Ressourcen und Schutz für eine Veränderung sichert.
  \item[Sustainment] Mechanismen, die eine Veränderung über Zeit verankern: Standards, Routinen, Messung, Coaching.
  \item[Trade-off-Protokoll] Dokumentation einer Entscheidung mit verworfenen Alternativen und Frühwarnsignal.
  \item[Truth Signal] Indikator, der zeigt, ob eine Veränderung wirklich angekommen ist (Adoption, Outcome, Pulse, Schattenprozess).
  \item[WIP-Limit (Change)] Obergrenze paralleler Change-Initiativen; verhindert Überlastung.
\end{description}

\vspace{1cm}
\noindent{\color{lpAccent}\rule{\linewidth}{0.6pt}}\\[4pt]
{\footnotesize\sffamily\color{lpGray}Lernplatt \textperiodcentered{} by Can Siebert \textperiodcentered{} Kurs 71 (Dig 42) \textperiodcentered{} Tag 15 \textperiodcentered{} \textcopyright{} 2026}

\end{document}
